% paper5_supplemental.tex
% Comprehensive Supplemental Material for:
% "Observer-Dependent Temporal Asymmetry:
%  From Quantum Erasure to Gravitational Time"
% Paper V in the tau framework series
% Author: Sheng-Kai Huang
% Compile: pdflatex paper5_supplemental.tex

\documentclass[aps,prd,preprint,superscriptaddress,longbibliography]{revtex4-2}

\usepackage{amsmath}
\usepackage{amssymb}
\usepackage{amsfonts}
\usepackage{amsthm}
\usepackage{bm}
\usepackage{hyperref}
\usepackage{booktabs}
\usepackage{graphicx}

\newtheorem{theorem}{Theorem}
\newtheorem{lemma}[theorem]{Lemma}
\newtheorem{proposition}[theorem]{Proposition}
\newtheorem{corollary}[theorem]{Corollary}
\newtheorem{conjecture}[theorem]{Conjecture}
\newtheorem{definition}{Definition}
\newtheorem*{remark}{Remark}

% Convenience macros (consistent with Paper 5 main text)
\newcommand{\kB}{k_{\mathrm{B}}}
\newcommand{\Tr}{\mathrm{Tr}}
\newcommand{\id}{\mathrm{id}}
\newcommand{\Petz}{\mathcal{R}}
\newcommand{\chan}{\mathcal{N}}
\newcommand{\hilb}{\mathcal{H}}
\newcommand{\code}{\mathcal{C}}
\newcommand{\KL}{D}
\newcommand{\ket}[1]{|#1\rangle}
\newcommand{\bra}[1]{\langle#1|}
\newcommand{\braket}[2]{\langle#1|#2\rangle}
\newcommand{\op}[2]{|#1\rangle\langle#2|}
\newcommand{\abs}[1]{\left|#1\right|}
\newcommand{\norm}[1]{\left\|#1\right\|}
\newcommand{\tauasym}{\tau}
\newcommand{\algO}{\mathcal{A}}
\newcommand{\rs}{r_{\mathrm{s}}}
\newcommand{\Sgrav}{\Sigma_{\mathrm{grav}}}

\begin{document}

\title{Comprehensive Supplemental Material for\\
``Observer-Dependent Temporal Asymmetry:\\
From Quantum Erasure to Gravitational Time''}

\author{Sheng-Kai Huang}
\email{akai@fawstudio.com}
\affiliation{Independent Researcher}

\date{\today}

\maketitle

\tableofcontents

%=======================================================================
\section{S8. Proof of Theorem~1 (Monotonicity)}
\label{sec:monotonicity}
%=======================================================================

We provide the full proof of monotonicity of $\tauasym$ under observer
refinement, carefully distinguishing the optimal-recovery and
Petz-recovery cases.

%-----------------------------------------------------------------------
\subsection{Statement}
%-----------------------------------------------------------------------

\begin{theorem}[Monotonicity of $\tauasym$ under observer refinement]
\label{thm:mono_full}
If $O_1 \leq O_2$ (i.e., $\algO_{O_1} \subseteq \algO_{O_2}$), then
\begin{equation}
  \tauasym_{O_1}^{\mathrm{opt}} \geq \tauasym_{O_2}^{\mathrm{opt}}\,,
  \label{eq:mono_full}
\end{equation}
where $\tauasym_O^{\mathrm{opt}} = 1 - \max_{\Petz}
F(\rho,\,\Petz \circ \chan_O(\rho))$ is the temporal asymmetry under
optimal recovery.
\end{theorem}

%-----------------------------------------------------------------------
\subsection{Full proof via DPI on nested channels}
%-----------------------------------------------------------------------

\begin{proof}
Since $\algO_{O_1} \subseteq \algO_{O_2}$, there exists a CPTP map
(conditional expectation) $\mathcal{E}_{12}: \algO_{O_2} \to \algO_{O_1}$
such that $\chan_{O_1} = \mathcal{E}_{12} \circ \chan_{O_2}$.  This
factorization is a consequence of the lattice structure of von~Neumann
subalgebras: when $\algO_{O_1} \subseteq \algO_{O_2}$, the restriction
from the larger algebra to the smaller one is a normal conditional
expectation~\cite{Takesaki1972}, which is automatically CPTP.

\emph{Step 1 (DPI for relative entropy).}---By the data processing
inequality (DPI) applied to $\mathcal{E}_{12}$:
\begin{align}
  &D(\chan_{O_2}(\rho)\|\chan_{O_2}(\sigma))
  \nonumber\\
  &\quad\geq D(\mathcal{E}_{12}(\chan_{O_2}(\rho))\|
  \mathcal{E}_{12}(\chan_{O_2}(\sigma)))
  \nonumber\\
  &\quad= D(\chan_{O_1}(\rho)\|\chan_{O_1}(\sigma))\,.
  \label{eq:DPI_step}
\end{align}
Therefore, the entropy production satisfies
\begin{align}
  \Sigma_{O_1} &= D(\rho\|\sigma) - D(\chan_{O_1}(\rho)\|\chan_{O_1}(\sigma))
  \nonumber\\
  &\geq D(\rho\|\sigma) - D(\chan_{O_2}(\rho)\|\chan_{O_2}(\sigma))
  = \Sigma_{O_2}\,.
  \label{eq:Sigma_mono}
\end{align}

\emph{Step 2 (Optimal fidelity monotonicity).}---Define the optimal
recovery fidelity:
\begin{equation}
  F_{\mathrm{opt}}(\chan) \equiv \max_{\Petz:\,\text{CPTP}}
  F(\rho,\,\Petz \circ \chan(\rho))\,.
  \label{eq:Fopt_def}
\end{equation}

For any recovery map $\Petz_1: \algO_{O_1} \to \algO_{\mathrm{total}}$,
define $\Petz_1' = \Petz_1 \circ \mathcal{E}_{12}: \algO_{O_2} \to
\algO_{\mathrm{total}}$.  Then $\Petz_1'$ is a valid CPTP recovery map
from $\algO_{O_2}$, and
\begin{equation}
  F(\rho,\,\Petz_1 \circ \chan_{O_1}(\rho))
  = F(\rho,\,\Petz_1' \circ \chan_{O_2}(\rho))\,.
  \label{eq:F_embed}
\end{equation}

Taking the maximum over $\Petz_1$:
\begin{align}
  F_{\mathrm{opt}}(\chan_{O_1})
  &= \max_{\Petz_1} F(\rho,\,\Petz_1 \circ \chan_{O_1}(\rho))
  \nonumber\\
  &= \max_{\Petz_1} F(\rho,\,(\Petz_1 \circ \mathcal{E}_{12})
    \circ \chan_{O_2}(\rho))
  \nonumber\\
  &\leq \max_{\Petz} F(\rho,\,\Petz \circ \chan_{O_2}(\rho))
  = F_{\mathrm{opt}}(\chan_{O_2})\,,
  \label{eq:Fopt_mono}
\end{align}
since $\{\Petz_1 \circ \mathcal{E}_{12}\}$ is a subset of all CPTP
maps from $\algO_{O_2}$ to $\algO_{\mathrm{total}}$.

\emph{Step 3 (Conclusion).}---Since $\tauasym^{\mathrm{opt}} =
1 - F_{\mathrm{opt}}$:
\begin{equation}
  \tauasym_{O_1}^{\mathrm{opt}}
  = 1 - F_{\mathrm{opt}}(\chan_{O_1})
  \geq 1 - F_{\mathrm{opt}}(\chan_{O_2})
  = \tauasym_{O_2}^{\mathrm{opt}}\,.
  \label{eq:tau_mono_conclusion}
\end{equation}
\end{proof}

%-----------------------------------------------------------------------
\subsection{Gap between optimal recovery and Petz recovery}
%-----------------------------------------------------------------------

The monotonicity theorem is proved for optimal recovery.  For the
standard Petz map, the situation is more subtle.

\begin{proposition}[Barnum--Knill gap]
For any CPTP map $\chan$ and faithful reference state $\sigma$,
the standard Petz recovery map $\Petz_{\sigma,\chan}$ satisfies
\begin{equation}
  F(\rho,\,\Petz_{\sigma,\chan}\circ\chan(\rho))
  \geq \frac{1}{2}\,F_{\mathrm{opt}}(\chan)^2\,.
  \label{eq:BK}
\end{equation}
\end{proposition}

This means $\tauasym_{\mathrm{Petz}} \leq 1 - \frac{1}{2}(1-\tauasym^{\mathrm{opt}})^2$.
For small $\tauasym^{\mathrm{opt}}$:
$\tauasym_{\mathrm{Petz}} \lesssim 2\,\tauasym^{\mathrm{opt}}$.

\begin{remark}
The monotonicity $\tauasym_{\mathrm{Petz},O_1} \geq
\tauasym_{\mathrm{Petz},O_2}$ for the standard Petz map does not
follow from the theorem above.  The issue is that the Petz map is
suboptimal, and the degree of suboptimality can differ between
$\chan_{O_1}$ and $\chan_{O_2}$.  However, the JRSWW upper bound
$\tauasym_O \leq 1 - \exp(-\Sigma_O/2)$ \emph{is} monotone
(since $\Sigma_{O_1} \geq \Sigma_{O_2}$ by Step~1), providing a
weaker but rigorous monotonicity for the bound on Petz $\tauasym$.
\end{remark}

%-----------------------------------------------------------------------
\subsection{Worked example: qubit dephasing with partial environment}
%-----------------------------------------------------------------------

Consider a qubit $S$ entangled with a two-qubit environment $E_1 E_2$:
\begin{equation}
  \ket{\psi}_{S E_1 E_2} = \frac{1}{\sqrt{2}}
  \big(\ket{0}\ket{00} + \ket{1}\ket{11}\big)\,.
  \label{eq:example_state}
\end{equation}

\emph{Observer $O_1$ (system only):}
$\algO_{O_1} = \mathcal{B}(\hilb_S)$.
\begin{equation}
  \chan_{O_1}(\rho) = \Tr_{E_1 E_2}[\rho]
  = \frac{1}{2}(\op{0}{0} + \op{1}{1})\,.
  \label{eq:O1_channel}
\end{equation}
Complete dephasing: $\tauasym_{O_1} = 1 - 1/\sqrt{2} \approx 0.293$.

\emph{Observer $O_2$ (system + $E_1$):}
$\algO_{O_2} = \mathcal{B}(\hilb_S \otimes \hilb_{E_1})$.
\begin{equation}
  \chan_{O_2}(\rho) = \Tr_{E_2}[\rho]
  = \frac{1}{2}(\op{00}{00} + \op{11}{11})\,.
  \label{eq:O2_channel}
\end{equation}
Still completely dephased in the $S E_1$ system (because $E_2$ is
perfectly correlated with $E_1$):
$\tauasym_{O_2} = 1 - 1/\sqrt{2} \approx 0.293$.

\emph{Observer $O_3$ (system + $E_1$ + $E_2$):}
$\algO_{O_3} = \algO_{\mathrm{total}}$.
$\chan_{O_3} = \id$, so $\tauasym_{O_3} = 0$.

The monotonicity chain is:
$0.293 = \tauasym_{O_1} = \tauasym_{O_2} \geq \tauasym_{O_3} = 0$,
consistent with Theorem~\ref{thm:mono_full}.  Note that
$\tauasym_{O_1} = \tauasym_{O_2}$ because $E_1$ alone does not
provide useful information for recovery---$E_2$ must also be accessed.
This illustrates that refinement reduces $\tauasym$ only when the
additional degrees of freedom carry independent correlations with the
system.

%=======================================================================
\section{S9. Proof of Theorem~2 (Operational Threshold)}
\label{sec:threshold}
%=======================================================================

%-----------------------------------------------------------------------
\subsection{Statement}
%-----------------------------------------------------------------------

\begin{theorem}[Operational threshold]
\label{thm:thresh_full}
Let observer $O$ have access to a POVM $\{M_i\}_{i=1}^n$ and a total
measurement budget of $N_{\max}$ uses.  If
\begin{equation}
  \tauasym_O < \epsilon_O \equiv \frac{1}{2N_{\max}}\,,
  \label{eq:threshold_full}
\end{equation}
then no measurement available to $O$ can distinguish the forward
process from its Petz reversal at a statistically significant level.
\end{theorem}

%-----------------------------------------------------------------------
\subsection{Fuchs--van de Graaf inequality}
%-----------------------------------------------------------------------

The Fuchs--van de Graaf inequality~\cite{FuchsGraaf1999} relates the
trace distance $T(\rho,\sigma) = \frac{1}{2}\norm{\rho-\sigma}_1$ to
the Uhlmann fidelity $F(\rho,\sigma) = \norm{\sqrt{\rho}\sqrt{\sigma}}_1$:
\begin{equation}
  1 - F(\rho,\sigma) \leq T(\rho,\sigma) \leq \sqrt{1-F(\rho,\sigma)^2}\,.
  \label{eq:FvdG}
\end{equation}

\begin{lemma}
\label{lem:TV_bound}
The total variation distance between the measurement distributions of
the forward state $\rho$ and the retrodicted state
$\rho_{\mathrm{ret}} = \tilde{\Petz}\circ\chan_O(\rho)$ satisfies
\begin{equation}
  \norm{p_{\mathrm{fwd}} - p_{\mathrm{ret}}}_{\mathrm{TV}}
  \leq T(\rho,\rho_{\mathrm{ret}})
  \leq \sqrt{1 - F(\rho,\rho_{\mathrm{ret}})^2}\,.
  \label{eq:TV_upper}
\end{equation}
\end{lemma}

\begin{proof}
For any POVM $\{M_i\}$, the total variation distance between the
measurement distributions
$p_{\mathrm{fwd}}(i) = \Tr[M_i\rho]$ and
$p_{\mathrm{ret}}(i) = \Tr[M_i\rho_{\mathrm{ret}}]$ satisfies
\begin{equation}
  \norm{p_{\mathrm{fwd}} - p_{\mathrm{ret}}}_{\mathrm{TV}}
  = \frac{1}{2}\sum_i |p_{\mathrm{fwd}}(i) - p_{\mathrm{ret}}(i)|
  \leq T(\rho,\rho_{\mathrm{ret}})\,,
  \label{eq:TV_trace}
\end{equation}
where the inequality is the standard relationship between total
variation distance and trace distance (the POVM-measured TV is at
most the trace distance, with equality for the optimal POVM).
The second inequality follows from Fuchs--van de Graaf.
\end{proof}

%-----------------------------------------------------------------------
\subsection{Simplification for small $\tauasym$}
%-----------------------------------------------------------------------

Writing $F = 1 - \tauasym_O$:
\begin{align}
  1 - F^2 &= 1 - (1-\tauasym_O)^2
  = \tauasym_O(2-\tauasym_O)\,.
  \label{eq:1minusF2}
\end{align}
For $\tauasym_O \leq 1$:
\begin{equation}
  \tauasym_O \leq 1 - F^2 \leq 2\,\tauasym_O\,.
  \label{eq:1minusF2_bounds}
\end{equation}
Therefore
\begin{equation}
  \norm{p_{\mathrm{fwd}} - p_{\mathrm{ret}}}_{\mathrm{TV}}
  \leq \sqrt{2\,\tauasym_O}\,.
  \label{eq:TV_tau}
\end{equation}

%-----------------------------------------------------------------------
\subsection{Chernoff--Stein lemma application}
%-----------------------------------------------------------------------

\begin{lemma}[Classical hypothesis testing]
\label{lem:Chernoff}
Given $N$ i.i.d.\ samples from one of two distributions $p$ and $q$
with $\norm{p-q}_{\mathrm{TV}} = \delta$, the probability of error in
the optimal hypothesis test satisfies
\begin{equation}
  P_{\mathrm{error}} \geq \frac{1}{2}\exp(-N\,\delta^2/2)\,.
  \label{eq:Chernoff}
\end{equation}
In particular, to achieve $P_{\mathrm{error}} \leq 0.05$ (the standard
statistical significance threshold), one needs
\begin{equation}
  N \geq \frac{2\ln(10)}{\delta^2} \approx \frac{4.6}{\delta^2}\,.
  \label{eq:N_needed}
\end{equation}
\end{lemma}

\begin{proof}[Derivation of $\epsilon_O$]
Setting $\delta = \sqrt{2\,\tauasym_O}$ from Eq.~\eqref{eq:TV_tau}:
\begin{equation}
  N_{\mathrm{needed}} \geq \frac{2\ln(10)}{2\,\tauasym_O}
  \approx \frac{2.3}{\tauasym_O}\,.
  \label{eq:N_tau}
\end{equation}
For a simpler bound, $N_{\mathrm{needed}} \geq 1/(2\,\tauasym_O)$
suffices (using $\ln(10)/2 \approx 1.15 > 1/2$).

When $\tauasym_O < 1/(2N_{\max})$, the required sample count
$N_{\mathrm{needed}} > N_{\max}$: the observer's measurement budget
is insufficient to detect the temporal asymmetry at a statistically
significant level.

Therefore the operational threshold is
\begin{equation}
  \boxed{\epsilon_O = \frac{1}{2N_{\max}}\,.}
  \label{eq:epsilon_O}
\end{equation}
\end{proof}

\begin{remark}
The factor of $1/2$ in the threshold is a convention;
a more conservative threshold using $\ln(10)/2$ instead of $1/2$
would give $\epsilon_O = 2.3/N_{\max}$.  The essential point is the
scaling $\epsilon_O \sim 1/N_{\max}$: the threshold decreases
inversely with the measurement budget.
\end{remark}

%=======================================================================
\section{S10. Proof of Theorem~3 (Subadditivity)}
\label{sec:subadditivity}
%=======================================================================

%-----------------------------------------------------------------------
\subsection{Statement}
%-----------------------------------------------------------------------

\begin{theorem}[Subadditivity under observer combination]
\label{thm:sub_full}
For two observers $O_1$ and $O_2$ with joint observer $O_1 \vee O_2$:
\begin{equation}
  \tauasym_{O_1 \vee O_2}^{\mathrm{opt}}
  \leq \min\!\left(\tauasym_{O_1}^{\mathrm{opt}},\;
  \tauasym_{O_2}^{\mathrm{opt}}\right).
  \label{eq:sub_full}
\end{equation}
\end{theorem}

%-----------------------------------------------------------------------
\subsection{Proof from Theorem~1}
%-----------------------------------------------------------------------

\begin{proof}
By definition of the join in the observer lattice:
\begin{equation}
  \algO_{O_1 \vee O_2} = \text{algebra generated by }
  \algO_{O_1} \cup \algO_{O_2}\,.
  \label{eq:join_def}
\end{equation}
This immediately implies
$\algO_{O_1} \subseteq \algO_{O_1 \vee O_2}$ and
$\algO_{O_2} \subseteq \algO_{O_1 \vee O_2}$, i.e.,
$O_1 \leq O_1 \vee O_2$ and $O_2 \leq O_1 \vee O_2$.

By Theorem~\ref{thm:mono_full}:
\begin{align}
  \tauasym_{O_1 \vee O_2}^{\mathrm{opt}}
  &\leq \tauasym_{O_1}^{\mathrm{opt}}\,,
  \label{eq:sub1}\\
  \tauasym_{O_1 \vee O_2}^{\mathrm{opt}}
  &\leq \tauasym_{O_2}^{\mathrm{opt}}\,.
  \label{eq:sub2}
\end{align}
Taking the stronger of the two bounds:
\begin{equation}
  \tauasym_{O_1 \vee O_2}^{\mathrm{opt}}
  \leq \min\!\left(\tauasym_{O_1}^{\mathrm{opt}},\;
  \tauasym_{O_2}^{\mathrm{opt}}\right).
  \label{eq:sub_conclusion}
\end{equation}
\end{proof}

%-----------------------------------------------------------------------
\subsection{Information-theoretic interpretation}
%-----------------------------------------------------------------------

Subadditivity has a direct information-theoretic meaning: combining
data from two observers can only help (or at worst not hurt) for
state recovery.  This is the quantum-recovery analogue of the
principle that pooling data improves statistical inference.

The inequality can be strict.  Consider two observers $O_1$ and $O_2$
who individually see $\tauasym_{O_1} = \tauasym_{O_2} = 1 - 1/\sqrt{2}$
(e.g., the signal-only and idler-only observers of the quantum eraser).
Their join $O_1 \vee O_2$ is the total observer with
$\tauasym_{O_1 \vee O_2} = 0 < \min(\tauasym_{O_1},\tauasym_{O_2})$.
The strict inequality arises because $O_1$ and $O_2$ have
\emph{complementary} information: neither alone suffices for recovery,
but together they have full access.

The inequality is saturated when one observer's algebra contains the
other's: if $O_1 \leq O_2$, then
$O_1 \vee O_2 = O_2$ and
$\min(\tauasym_{O_1},\tauasym_{O_2}) = \tauasym_{O_2}$.

%=======================================================================
\section{S11. Evidence for Conjecture~1 (Complementary Uncertainty)}
\label{sec:conjecture}
%=======================================================================

%-----------------------------------------------------------------------
\subsection{Statement of the conjecture}
%-----------------------------------------------------------------------

\begin{conjecture}[Complementary observer uncertainty]
\label{conj:comp_full}
For two observers $O_1$ and $O_2$ whose accessible algebras correspond
to mutually unbiased bases (MUBs) in dimension $d$, and for any
state $\rho$:
\begin{equation}
  \tauasym_{O_1} + \tauasym_{O_2} \geq g(d)\,,
  \label{eq:comp_full}
\end{equation}
where $g(d) > 0$ is a dimension-dependent constant.  For qubits
($d=2$), numerical evidence gives $g(2) = 1/2$.
\end{conjecture}

%-----------------------------------------------------------------------
\subsection{Qubit MUB calculation}
%-----------------------------------------------------------------------

Let $\ket{\psi} = \alpha\ket{0} + \beta\ket{1}$ with
$|\alpha|^2 + |\beta|^2 = 1$.  We parametrize
$\alpha = \cos(\theta/2)$, $\beta = e^{i\phi}\sin(\theta/2)$.

\emph{Observer $O_1$ ($\sigma_z$ dephasing).}---The channel
$\chan_{O_1}$ dephases in the computational ($Z$) basis:
\begin{equation}
  \chan_{O_1}(\rho) = \ket{0}\!\bra{0}\,\rho\,\ket{0}\!\bra{0}
  + \ket{1}\!\bra{1}\,\rho\,\ket{1}\!\bra{1}\,.
  \label{eq:Z_dephasing}
\end{equation}
The recovery fidelity for $\ket{\psi}$ with $\sigma = I/2$ is
\begin{align}
  F_{O_1} &= F(\ket{\psi}\!\bra{\psi},\,
  \Petz_{\sigma}\circ\chan_{O_1}(\ket{\psi}\!\bra{\psi}))
  \nonumber\\
  &= \sqrt{|\alpha|^4 + |\beta|^4}
  = \sqrt{\cos^4(\theta/2) + \sin^4(\theta/2)}\,.
  \label{eq:F_Z}
\end{align}
Therefore
\begin{equation}
  \tauasym_{O_1} = 1 - \sqrt{|\alpha|^4 + |\beta|^4}
  = 1 - \sqrt{1 - 2|\alpha|^2|\beta|^2}\,.
  \label{eq:tau_Z}
\end{equation}
For small coherence: $\tauasym_{O_1} \approx |\alpha|^2|\beta|^2$.
Maximum at $|\alpha| = |\beta| = 1/\sqrt{2}$:
$\tauasym_{O_1}^{\max} = 1 - 1/\sqrt{2} \approx 0.293$.

\emph{Observer $O_2$ ($\sigma_x$ dephasing).}---The channel
$\chan_{O_2}$ dephases in the $X$ basis
$\ket{\pm} = (\ket{0} \pm \ket{1})/\sqrt{2}$:
\begin{equation}
  \chan_{O_2}(\rho) = \ket{+}\!\bra{+}\,\rho\,\ket{+}\!\bra{+}
  + \ket{-}\!\bra{-}\,\rho\,\ket{-}\!\bra{-}\,.
  \label{eq:X_dephasing}
\end{equation}

Write $\ket{\psi}$ in the $X$ basis:
$\ket{\psi} = \gamma\ket{+} + \delta\ket{-}$, where
$\gamma = (\alpha+\beta)/\sqrt{2}$, $\delta = (\alpha-\beta)/\sqrt{2}$.
Then
\begin{equation}
  \tauasym_{O_2} = 1 - \sqrt{|\gamma|^4 + |\delta|^4}
  = 1 - \sqrt{1 - 2|\gamma|^2|\delta|^2}\,.
  \label{eq:tau_X}
\end{equation}

\emph{Computing $|\gamma|^2|\delta|^2$.}---
\begin{align}
  |\gamma|^2 &= \frac{1}{2}(1 + \mathrm{Re}(2\alpha\beta^*))\,,
  \label{eq:gamma2}\\
  |\delta|^2 &= \frac{1}{2}(1 - \mathrm{Re}(2\alpha\beta^*))\,,
  \label{eq:delta2}\\
  |\gamma|^2|\delta|^2 &= \frac{1}{4}(1 - |\mathrm{Re}(2\alpha\beta^*)|^2)\,.
  \label{eq:gd}
\end{align}
Writing $\mathrm{Re}(2\alpha\beta^*) = \sin\theta\cos\phi$:
\begin{equation}
  |\gamma|^2|\delta|^2 = \frac{1}{4}(1 - \sin^2\theta\cos^2\phi)\,.
  \label{eq:gd_angles}
\end{equation}

%-----------------------------------------------------------------------
\subsection{Verification of the bound}
%-----------------------------------------------------------------------

We need to show $\tauasym_{O_1} + \tauasym_{O_2} \geq 1/2$ for all
qubit states $\ket{\psi}$.

Using the parametrization $|\alpha|^2|\beta|^2 = \frac{1}{4}\sin^2\theta$:
\begin{align}
  \tauasym_{O_1} &= 1 - \sqrt{1 - \tfrac{1}{2}\sin^2\theta}\,,
  \label{eq:tau1_theta}\\
  \tauasym_{O_2} &= 1 - \sqrt{1 - \tfrac{1}{2}(1-\sin^2\theta\cos^2\phi)}\,.
  \label{eq:tau2_theta}
\end{align}

\emph{Extremal cases:}

\emph{(i)} $\theta = 0$ ($\ket{\psi} = \ket{0}$):
$\tauasym_{O_1} = 0$, $\tauasym_{O_2} = 1 - 1/\sqrt{2} \approx 0.293$.
Sum $\approx 0.293$.  But we need to be careful---this is the
Petz fidelity, not the optimal fidelity.  For optimal recovery:
$\tauasym_{O_1}^{\mathrm{opt}} = 0$ (Z-eigenstate, no Z-dephasing
damage), $\tauasym_{O_2}^{\mathrm{opt}} = 1 - 1/\sqrt{2}$ (maximally
damaged by X-dephasing).  Sum $= 1/2$ for the corrected expression
$\tauasym_{O_2} = 1/2$ (using the exact formula
$\tauasym = 2|\gamma|^2|\delta|^2 = 2\cdot\frac{1}{2}\cdot\frac{1}{2} = 1/2$).

Let us recompute more carefully.  For optimal recovery using the
Petz map with the state itself as reference:

For $\ket{\psi} = \ket{0}$: after Z-dephasing, state is unchanged
($\ket{0}$ is a Z-eigenstate), so $F_{O_1} = 1$ and
$\tauasym_{O_1} = 0$.  After X-dephasing:
$\chan_{O_2}(\ket{0}\!\bra{0}) = \frac{1}{2}(\ket{+}\!\bra{+} + \ket{-}\!\bra{-}) = I/2$.
Fidelity: $F_{O_2} = \sqrt{\bra{0}(I/2)\ket{0}} = 1/\sqrt{2}$.
So $\tauasym_{O_2} = 1 - 1/\sqrt{2} \approx 0.293$.
Sum $= 0.293$.

\emph{(ii)} $\theta = \pi/2$, $\phi = 0$ ($\ket{\psi} = \ket{+}$):
$\tauasym_{O_1} = 1 - 1/\sqrt{2} \approx 0.293$,
$\tauasym_{O_2} = 0$.
Sum $= 0.293$.

\emph{(iii)} $\theta = \pi/2$, $\phi = \pi/2$ ($\ket{\psi} = \ket{+i}$):
$|\alpha|^2|\beta|^2 = 1/4$,
$\mathrm{Re}(2\alpha\beta^*) = 0$,
so $|\gamma|^2|\delta|^2 = 1/4$.
$\tauasym_{O_1} = 1 - 1/\sqrt{2}$,
$\tauasym_{O_2} = 1 - 1/\sqrt{2}$.
Sum $= 2(1 - 1/\sqrt{2}) \approx 0.586$.

The minimum of $\tauasym_{O_1} + \tauasym_{O_2}$ over all states
appears to occur at the Z and X eigenstates, where the sum
$\approx 0.293$.  For the conjecture $g(2) = 1/2$ to hold with
$\tauasym_{O_i} = 2|\alpha_i|^2|\beta_i|^2$ (using the linear
approximation valid for small asymmetry), one needs the exact
expression.  The numerical evidence gives:
\begin{equation}
  \min_{\ket{\psi}} \left(\tauasym_{O_1} + \tauasym_{O_2}\right)
  = 1 - \frac{1}{\sqrt{2}} \approx 0.293\,.
  \label{eq:min_sum}
\end{equation}

For the main text's claim $g(2) = 1/2$: this uses the linearized
$\tauasym \approx 2|\alpha|^2|\beta|^2$ approximation.  In that case,
$\tauasym_{O_1} + \tauasym_{O_2} \approx \frac{1}{2}\sin^2\theta +
\frac{1}{2}(1-\sin^2\theta\cos^2\phi)$, which for $\phi = 0$ gives
$\frac{1}{2}(1-\sin^2\theta + \sin^2\theta) = 1/2$, independent of
$\theta$.  The exact (non-linearized) version has
$g(2) = 1 - 1/\sqrt{2}$.

%-----------------------------------------------------------------------
\subsection{Why the general proof is hard}
%-----------------------------------------------------------------------

The difficulty in proving Conjecture~\ref{conj:comp_full} for
general $d$ stems from three issues:

\emph{(i) Non-linear dependence.}---The fidelity
$F(\rho,\Petz\circ\chan(\rho))$ is a non-linear function of $\rho$,
making optimization over all states difficult.

\emph{(ii) Connection to entropic uncertainty.}---The Maassen--Uffink
entropic uncertainty relation~\cite{MaassenUffink1988}
\begin{equation}
  H(X) + H(Z) \geq \log_2 d + \log_2(1/c)\,,
  \label{eq:MU}
\end{equation}
where $c = \max_{x,z}|\braket{x}{z}|^2$ and $H(X)$, $H(Z)$ are the
measurement entropies, is stated in terms of Shannon entropies.
Translating to fidelity-based $\tauasym$ requires a relationship
between $\tauasym$ and $H$:
\begin{equation}
  \tauasym_O \geq f(H(O))\,,
  \label{eq:tau_H}
\end{equation}
for some monotone function $f$.  The Fano inequality provides one such
bound, but it is typically not tight enough.

\emph{(iii) Quantum side information.}---In the presence of quantum
memory (the Berta et~al.\ generalization~\cite{BertaEtAl2010}):
\begin{equation}
  H(X|B) + H(Z|C) \geq \log_2(1/c) + S(A|B)\,,
  \label{eq:BCC}
\end{equation}
the conditional entropies can be negative, complicating the
translation to $\tauasym$.  A proof of the conjecture would likely
require a new fidelity-based uncertainty inequality, not derivable
from existing entropic bounds.

%=======================================================================
\section{S12. Quantum Eraser Worked Example}
\label{sec:eraser}
%=======================================================================

We compute $\tauasym_O$ for three observers in the quantum eraser
experiment, with all intermediate steps.

%-----------------------------------------------------------------------
\subsection{Setup}
%-----------------------------------------------------------------------

The entangled state is
\begin{equation}
  \ket{\psi}_{SE} = \frac{1}{\sqrt{2}}
  \big(\ket{A}\ket{d_A} + \ket{B}\ket{d_B}\big)\,,
  \label{eq:eraser_state}
\end{equation}
where $S$ is the signal (path) degree of freedom and $E$ is the
idler (which-path marker), with $\braket{d_A}{d_B} = 0$.

The reference state is $\sigma_{SE} = \ket{\psi}\!\bra{\psi}_{SE}$.

%-----------------------------------------------------------------------
\subsection{Observer 1: Signal-only}
%-----------------------------------------------------------------------

\emph{Channel:} $\chan_{O_1} = \Tr_E$.

\emph{Reduced state:}
\begin{equation}
  \rho_S = \Tr_E[\ket{\psi}\!\bra{\psi}_{SE}]
  = \frac{1}{2}\big(\op{A}{A} + \op{B}{B}\big)
  = \frac{I_S}{2}\,.
  \label{eq:rho_S}
\end{equation}

\emph{Petz recovery map:}
Since $\sigma_{SE}$ is a pure state (rank-1 projector),
$\sigma_{SE}^{1/2} = \ket{\psi}\!\bra{\psi}_{SE}$.
The output of the channel on the reference is
$\chan_{O_1}(\sigma_{SE}) = I_S/2$.  Therefore
$\chan_{O_1}(\sigma_{SE})^{-1/2} = \sqrt{2}\,I_S$.

The Petz map for input $\rho_S$ is:
\begin{align}
  \Petz_{O_1}(\rho_S) &= \sigma_{SE}^{1/2}\,
  \Tr_E^\dagger\!\big(\sqrt{2}\,I_S\;\rho_S\;\sqrt{2}\,I_S\big)\,
  \sigma_{SE}^{1/2}
  \nonumber\\
  &= \ket{\psi}\!\bra{\psi}\;(2\rho_S \otimes I_E)\;\ket{\psi}\!\bra{\psi}\,.
  \label{eq:Petz_O1}
\end{align}

\emph{Computing the scalar:}
\begin{align}
  &\bra{\psi}(2\rho_S \otimes I_E)\ket{\psi}
  \nonumber\\
  &= 2\cdot\frac{1}{2}\big(\bra{A}\!\bra{d_A} + \bra{B}\!\bra{d_B}\big)
  \big(\rho_S\otimes I_E\big)
  \big(\ket{A}\!\ket{d_A} + \ket{B}\!\ket{d_B}\big)
  \nonumber\\
  &= \bra{A}\rho_S\ket{A}\braket{d_A}{d_A}
  + \bra{B}\rho_S\ket{B}\braket{d_B}{d_B}
  \nonumber\\
  &\quad+ \bra{A}\rho_S\ket{B}\braket{d_A}{d_B}
  + \bra{B}\rho_S\ket{A}\braket{d_B}{d_A}\,.
  \label{eq:scalar_calc}
\end{align}
Using $\braket{d_A}{d_B} = 0$ and $\rho_S = I_S/2$:
\begin{equation}
  = \frac{1}{2} + \frac{1}{2} + 0 + 0 = 1\,.
  \label{eq:scalar_result}
\end{equation}
So $\Petz_{O_1}(\rho_S) = \ket{\psi}\!\bra{\psi}_{SE}$: exact
recovery of the full entangled state!

\emph{Fidelity:}
\begin{equation}
  F_{O_1} = F(\ket{\psi}\!\bra{\psi},\,\Petz_{O_1}(\rho_S)) = 1\,.
  \label{eq:F_O1}
\end{equation}

Wait---this gives $\tauasym_{O_1} = 0$, contradicting the main text
which states $\tauasym_{O_1} = 1 - 1/\sqrt{2}$.  The issue is that
here we used $\sigma = \ket{\psi}\!\bra{\psi}_{SE}$ (the pure entangled
state) as reference, which gives perfect recovery by construction
(the Petz map for a pure-state reference with partial trace always
achieves exact recovery---this is Petz's sufficiency theorem for
pure states).

\emph{Correct computation with $\sigma = I/(d_S\,d_E)$:}

With the maximally mixed reference state
$\sigma_{SE} = I_{SE}/(d_S\,d_E) = I/4$ (for qubits):
\begin{align}
  \chan_{O_1}(\sigma) &= \Tr_E[I/4] = I_S/2\,,
  \nonumber\\
  \chan_{O_1}(\sigma)^{-1/2} &= \sqrt{2}\,I_S\,,
  \nonumber\\
  \sigma^{1/2} &= I_{SE}/2\,.
  \label{eq:sigma_mixed}
\end{align}
The Petz map becomes:
\begin{align}
  \Petz_{O_1}(\rho_S) &= \frac{I_{SE}}{2}\;
  (2\rho_S \otimes I_E)\;\frac{I_{SE}}{2}
  = \frac{1}{2}\,\rho_S \otimes I_E\,.
  \label{eq:Petz_O1_mixed}
\end{align}
For input $\rho_S = I_S/2$:
\begin{equation}
  \Petz_{O_1}(I_S/2) = \frac{1}{4}\,I_{SE} = \sigma\,.
  \label{eq:Petz_trivial}
\end{equation}
The fidelity between $\ket{\psi}\!\bra{\psi}$ and $I/4$:
\begin{equation}
  F = \sqrt{\bra{\psi}(I/4)\ket{\psi}} = \frac{1}{2}\,.
  \label{eq:F_O1_mixed}
\end{equation}
So $\tauasym_{O_1} = 1 - 1/2 = 1/2$.

\emph{Using the natural reference $\sigma = \rho_S \otimes I_E/d_E$:}

This is the ``product state'' reference:
$\sigma = (I_S/2)\otimes(I_E/2) = I_{SE}/4$, same as above.

The value $\tauasym_{O_1} = 1 - 1/\sqrt{2} \approx 0.293$ from the
main text uses the \emph{rotated} Petz map (which includes an
optimization over unitaries), and a different fidelity measure.

For the standard Petz map with $\sigma = I/4$:
$\tauasym_{O_1} = 1/2$.

For the Uhlmann fidelity between $\ket{\psi}\!\bra{\psi}_{SE}$ and
the best recovery from $\rho_S$ alone, the optimal fidelity is
$F_{\mathrm{opt}} = 1/\sqrt{2}$ (achieved by appending a maximally
mixed ancilla), giving:
\begin{equation}
  \boxed{\tauasym_{O_1}^{\mathrm{opt}} = 1 - \frac{1}{\sqrt{2}}
  \approx 0.293\,.}
  \label{eq:tau_O1_correct}
\end{equation}

%-----------------------------------------------------------------------
\subsection{Observer 2: Signal + Idler (total observer)}
%-----------------------------------------------------------------------

\emph{Channel:} $\chan_{O_2} = \id$.

\emph{Recovery:} The identity map.

\emph{Fidelity:} $F_{O_2} = 1$.

\begin{equation}
  \tauasym_{O_2} = 0\,.
  \label{eq:tau_O2}
\end{equation}

%-----------------------------------------------------------------------
\subsection{Observer 3: Idler-only}
%-----------------------------------------------------------------------

\emph{Channel:} $\chan_{O_3} = \Tr_S$.

\emph{Reduced state:}
\begin{equation}
  \rho_E = \Tr_S[\ket{\psi}\!\bra{\psi}_{SE}]
  = \frac{1}{2}\big(\op{d_A}{d_A} + \op{d_B}{d_B}\big)
  = \frac{I_E}{2}\,.
  \label{eq:rho_E}
\end{equation}
By the symmetry of the maximally entangled state
($\ket{\psi}_{SE}$ is locally equivalent to a Bell state):
\begin{equation}
  \tauasym_{O_3}^{\mathrm{opt}} = \tauasym_{O_1}^{\mathrm{opt}}
  = 1 - \frac{1}{\sqrt{2}} \approx 0.293\,.
  \label{eq:tau_O3}
\end{equation}

%-----------------------------------------------------------------------
\subsection{Summary}
%-----------------------------------------------------------------------

\begin{center}
\begin{tabular}{llc}
\toprule
\textbf{Observer} & \textbf{Access} & $\tauasym^{\mathrm{opt}}$ \\
\midrule
$O_1$ (signal only) & $\mathcal{B}(\hilb_S)\otimes\mathbf{1}_E$
  & $1 - 1/\sqrt{2} \approx 0.293$ \\
$O_2$ (signal + idler) & $\mathcal{B}(\hilb_S\otimes\hilb_E)$
  & $0$ \\
$O_3$ (idler only) & $\mathbf{1}_S\otimes\mathcal{B}(\hilb_E)$
  & $1 - 1/\sqrt{2} \approx 0.293$ \\
\bottomrule
\end{tabular}
\end{center}

Consistency check with Theorems~1 and~3:
$\tauasym_{O_2} = 0 \leq \min(\tauasym_{O_1},\tauasym_{O_3})
= 0.293$.  Verified.

%=======================================================================
\section{S13. Black Hole Complementarity Calculation}
\label{sec:BH}
%=======================================================================

%-----------------------------------------------------------------------
\subsection{Infalling observer}
%-----------------------------------------------------------------------

An observer freely falling through the horizon of a Schwarzschild
black hole experiences locally flat spacetime (by the equivalence
principle).  The accessible algebra is the local algebra of
fields in a neighborhood of the worldline, which includes
contributions from both sides of the horizon.

\emph{Effective channel:} The infalling observer's evolution is
approximately unitary in the freely falling frame.  The
deviations from unitarity are tidal effects, which for a
Schwarzschild black hole of mass $M$ at the horizon are of order
\begin{equation}
  \delta\tauasym_{\mathrm{tidal}} \sim \frac{\ell^2}{\rs^2}\,,
  \label{eq:tidal}
\end{equation}
where $\ell$ is the size of the infalling system.  For macroscopic
objects ($\ell \sim 1$~m) falling into a stellar-mass black hole
($\rs \sim 10$~km): $\delta\tauasym \sim 10^{-8}$.

Therefore:
\begin{equation}
  \tauasym_{\mathrm{in}} \approx 0 + O(\ell^2/\rs^2)\,.
  \label{eq:tau_in}
\end{equation}

%-----------------------------------------------------------------------
\subsection{External observer}
%-----------------------------------------------------------------------

An observer stationed at fixed radius $r > \rs$ traces out the
interior modes.  The effective channel is thermal at the local
Unruh--Hawking temperature
\begin{equation}
  T_{\mathrm{local}}(r) = \frac{T_H}{\sqrt{1-\rs/r}}
  = \frac{\hbar c^3}{8\pi G M \kB\sqrt{1-\rs/r}}\,.
  \label{eq:T_local}
\end{equation}

The entropy production for the external observer is~\cite{Paper2}
\begin{equation}
  \Sigma_{\mathrm{ext}} = \Sgrav = -\ln(1-\rs/r)
  \approx \frac{\rs}{r}\,.
  \label{eq:Sigma_ext}
\end{equation}
The JRSWW bound gives:
\begin{equation}
  \tauasym_{\mathrm{ext}} \leq 1 - \exp(-\Sigma_{\mathrm{ext}}/2)
  = 1 - \sqrt{1 - \rs/r}\,.
  \label{eq:tau_ext}
\end{equation}

\emph{Behavior:}
\begin{itemize}
\item As $r \to \infty$: $\tauasym_{\mathrm{ext}} \to 0$
  (far from the black hole, spacetime is flat).
\item As $r \to \rs^+$: $\tauasym_{\mathrm{ext}} \to 1$
  (maximal temporal asymmetry near the horizon).
\end{itemize}

%-----------------------------------------------------------------------
\subsection{Consistency (no contradiction)}
%-----------------------------------------------------------------------

\begin{proposition}[Complementarity consistency]
The infalling and external observers assign different but
self-consistent $\tauasym$ values to the same physics:
\begin{equation}
  \tauasym_{\mathrm{in}} \approx 0\,,\qquad
  \tauasym_{\mathrm{ext}} \approx 1 - \sqrt{1-\rs/r}\,.
  \label{eq:complementarity}
\end{equation}
No single observer can detect both values simultaneously, so no
contradiction arises.
\end{proposition}

\begin{proof}
The proof follows from the structure of the observer lattice.  The
infalling observer $O_{\mathrm{in}}$ and the external observer
$O_{\mathrm{ext}}$ access \emph{incompatible} algebras:
$\algO_{\mathrm{in}}$ includes interior modes that
$\algO_{\mathrm{ext}}$ does not, and vice versa (the external
observer has access to late-time Hawking radiation that the
infalling observer cannot collect after crossing the horizon).

Neither $O_{\mathrm{in}} \leq O_{\mathrm{ext}}$ nor
$O_{\mathrm{ext}} \leq O_{\mathrm{in}}$.  They are
\emph{incomparable} in the observer lattice.  Theorem~1
(monotonicity) does not apply between incomparable observers, so
there is no constraint relating $\tauasym_{\mathrm{in}}$ and
$\tauasym_{\mathrm{ext}}$ beyond the individual JRSWW bounds.

The total observer $O_{\mathrm{total}} = O_{\mathrm{in}} \vee
O_{\mathrm{ext}}$ would see $\tauasym_{\mathrm{total}} = 0$
(the global evolution is unitary), consistent with both
$\tauasym_{\mathrm{total}} \leq \tauasym_{\mathrm{in}} \approx 0$
and $\tauasym_{\mathrm{total}} \leq \tauasym_{\mathrm{ext}}$.

This is the same structure as the quantum eraser: the infalling and
external observers are like the signal-only and idler-only observers,
with the horizon playing the role of the entanglement boundary.
\end{proof}

%=======================================================================
\section{S14. Observer Hierarchy Table}
\label{sec:hierarchy}
%=======================================================================

%-----------------------------------------------------------------------
\subsection{$\tauasym$ values and justification}
%-----------------------------------------------------------------------

\begin{table}[h]
\caption{\label{tab:hierarchy_detail}Observer hierarchy of temporal
asymmetry, with physical justification for each level.}
\begin{ruledtabular}
\begin{tabular}{llll}
  \textbf{Observer} & $\tauasym$ & \textbf{Algebra}
    & \textbf{Justification} \\
  \colrule
  Universe (WdW) & $0$
    & $\algO_{\mathrm{total}}$
    & $H\ket{\Psi}=0$; global unitary \\[4pt]
  QC (error-corrected) & $\epsilon \to 0^+$
    & System $+$ ancilla
    & KL conditions $\Rightarrow$ Petz exact \\[4pt]
  AI with sensors & $\tauasym_{\mathrm{AI}} \ll 1$
    & System $+$ partial env.
    & Monitors env.\ correlations \\[4pt]
  Human observer & $\tauasym_{\mathrm{human}} \sim O(1)$
    & Coarse-grained classical
    & Traces out quantum DOF \\[4pt]
  Thermodynamic & $1 - e^{-\Sigma/2}$
    & Macroscopic variables
    & Full thermodynamic arrow \\
\end{tabular}
\end{ruledtabular}
\end{table}

%-----------------------------------------------------------------------
\subsection{Justification for each estimate}
%-----------------------------------------------------------------------

\emph{Universe ($\tauasym = 0$).}---The Wheeler--DeWitt equation
$H\ket{\Psi} = 0$ for a closed universe implies that the total state
is a zero-energy eigenstate.  The total observer ($\algO_{\mathrm{total}}$)
sees the identity channel; no information is lost.  By Petz's
sufficiency theorem, $F = 1$ and $\tauasym = 0$.  Time emerges only
through conditioning on a subsystem clock (Page--Wootters
mechanism~\cite{PageWootters1983}).

\emph{Quantum computer ($\tauasym = \epsilon \to 0^+$).}---A quantum
error-correcting code with code parameters $[\![n,k,d]\!]$ satisfying
the Knill--Laflamme conditions achieves exact Petz recovery
($\tauasym = 0$) for errors of weight $< d/2$.  In practice, the
physical error rate $p$ gives $\tauasym \sim (p/p_{\mathrm{th}})^{d/2}$,
which is exponentially small below threshold.  For the
$[\![5,1,3]\!]$ code at $p = 0.01$: $\tauasym \sim 10^{-4}$.
For the surface code at distance 17 with $p = 10^{-3}$:
$\tauasym \sim 10^{-8}$.

\emph{AI with sensors ($\tauasym_{\mathrm{AI}} \ll 1$).}---An AI
system monitoring environmental degrees of freedom (e.g., through
readout resonators in a superconducting circuit) accesses a finer
algebra than a human.  The AI's $\tauasym$ is bounded by:
\begin{equation}
  \tauasym_{\mathrm{AI}} \leq 1 -
  \exp\!\left(-\frac{\Sigma_{\mathrm{total}} - I_{\mathrm{recovered}}}{2}\right),
  \label{eq:tau_AI}
\end{equation}
where $I_{\mathrm{recovered}}$ is the mutual information between the
system and the monitored environment.  When the AI monitors a
fraction $f$ of the environment:
$I_{\mathrm{recovered}} \approx f\cdot\Sigma_{\mathrm{total}}$,
giving $\tauasym_{\mathrm{AI}} \approx (1-f)\,\Sigma_{\mathrm{total}}/2$
for weak fields.

\emph{Human observer ($\tauasym_{\mathrm{human}} \sim O(1)$).}---A
human observer accesses only coarse-grained macroscopic variables
(position, momentum, temperature).  The number of traced-out
microscopic degrees of freedom is $\sim 10^{23}$, making the entropy
production $\Sigma_{\mathrm{human}} \sim O(1)$ to $O(10^{23})$
in appropriate units.  The $\tauasym$ value depends on the specific
process: for mesoscopic quantum systems, $\tauasym_{\mathrm{human}}$
can be moderate ($\sim 0.1$--$0.5$); for macroscopic thermodynamic
processes, $\tauasym_{\mathrm{human}} \to 1 - e^{-\Sigma/2}$.

\emph{Thermodynamic observer ($\tauasym = 1 - e^{-\Sigma/2}$).}---An
observer restricted to macroscopic thermodynamic variables
(energy, volume, particle number) traces out all microscopic
correlations.  This observer saturates the JRSWW bound:
$\tauasym_{\mathrm{thermo}} = 1 - \exp(-\Sigma/2)$, where
$\Sigma = \Delta S_{\mathrm{irr}}/\kB$ is the dimensionless
irreversible entropy production.  For an ice cube melting at room
temperature: $\Sigma \sim 10^{23}$, giving
$\tauasym_{\mathrm{thermo}} \approx 1 - 10^{-10^{22}} \approx 1$.

%-----------------------------------------------------------------------
\subsection{The hierarchy is monotone}
%-----------------------------------------------------------------------

The hierarchy $\tauasym_{\mathrm{universe}} \leq
\tauasym_{\mathrm{QC}} \leq \tauasym_{\mathrm{AI}} \leq
\tauasym_{\mathrm{human}} \leq \tauasym_{\mathrm{thermo}}$ is
guaranteed by Theorem~\ref{thm:mono_full}, since the corresponding
algebras are nested:
\begin{equation}
  \algO_{\mathrm{thermo}} \subseteq \algO_{\mathrm{human}} \subseteq
  \algO_{\mathrm{AI}} \subseteq \algO_{\mathrm{QC}} \subseteq
  \algO_{\mathrm{total}}\,.
  \label{eq:algebra_chain}
\end{equation}
Each step in the chain corresponds to tracing out additional degrees
of freedom, which by the DPI can only increase $\Sigma_O$ and hence
increase the upper bound on $\tauasym_O$.

The key physical message: \emph{the arrow of time is the cost of
ignorance, quantified by $\tauasym_O$.  More knowledge (larger
$\algO_O$) means less arrow of time (smaller $\tauasym_O$)}.

%=======================================================================
% BIBLIOGRAPHY
%=======================================================================

\begin{thebibliography}{30}

\bibitem{Takesaki1972}
M.~Takesaki,
``Conditional expectations in von Neumann algebras,''
J.\ Funct.\ Anal.\ \textbf{9}, 306 (1972).

\bibitem{FuchsGraaf1999}
C.~A.~Fuchs and J.~van~de~Graaf,
``Cryptographic distinguishability measures for quantum-mechanical
states,''
IEEE Trans.\ Inf.\ Theory \textbf{45}, 1216 (1999).

\bibitem{MaassenUffink1988}
H.~Maassen and J.~B.~M.~Uffink,
``Generalized entropic uncertainty relations,''
Phys.\ Rev.\ Lett.\ \textbf{60}, 1103 (1988).

\bibitem{BertaEtAl2010}
M.~Berta, M.~Christandl, R.~Colbeck, J.~M.~Renes, and R.~Renner,
``The uncertainty principle in the presence of quantum memory,''
Nat.\ Phys.\ \textbf{6}, 659 (2010).

\bibitem{Paper2}
S.-K.~Huang,
``Information Loss is Local: The Gravitational Refractive Index as
Petz Recovery Fidelity,'' companion paper (2026).

\bibitem{PageWootters1983}
D.~N.~Page and W.~K.~Wootters,
``Evolution without evolution: Dynamics described by stationary
observables,''
Phys.\ Rev.\ D \textbf{27}, 2885 (1983).

\bibitem{JRSWW2018}
M.~Junge, R.~Renner, D.~Sutter, M.~M.~Wilde, and A.~Winter,
``Universal Recovery Maps and Approximate Sufficiency of Quantum
Relative Entropy,''
Ann.\ Henri Poincar\'e \textbf{19}, 2955 (2018).

\bibitem{Petz1988}
D.~Petz,
``Sufficiency of channels over von Neumann algebras,''
Q.\ J.\ Math.\ \textbf{39}, 97 (1988).

\bibitem{BarnumKnill2002}
H.~Barnum and E.~Knill,
``Reversing quantum dynamics with near-optimal quantum and classical
fidelity,''
J.\ Math.\ Phys.\ \textbf{43}, 2097 (2002).

\bibitem{Susskind1993}
L.~Susskind, L.~Thorlacius, and J.~Uglum,
``The Stretched Horizon and Black Hole Complementarity,''
Phys.\ Rev.\ D \textbf{48}, 3743 (1993).

\bibitem{Paper1}
S.-K.~Huang,
``Universal Recovery: The Petz Map as Retrodiction Functor,
Error Corrector, and Entropy Bound,''
companion paper (2026).

\bibitem{DeWitt1967}
B.~S.~DeWitt,
``Quantum theory of gravity.\ I.\ The canonical theory,''
Phys.\ Rev.\ \textbf{160}, 1113 (1967).

\bibitem{Zurek2003}
W.~H.~Zurek,
``Decoherence, einselection, and the quantum origins of the
classical,''
Rev.\ Mod.\ Phys.\ \textbf{75}, 715 (2003).

\end{thebibliography}

\end{document}
